% Chapter 14 -- An 8-bit ALU
\chapter{An 8-bit ALU}\label{ch:alu8}

Every processor you have ever used --- the CPU in your laptop, the
microcontroller on a sensor board, the soft core we will build in
ch16 --- has one block at its centre that answers questions like
``what is $a+b$?'', ``is $a$ bigger than $b$?'' and ``give me the bits
of $a$ that are also set in $b$.'' That block is the
\emph{arithmetic-logic unit}, and it is not magical: it is a small
stack of ordinary circuits --- an adder, some gates, two shifters and
a big mux --- sharing two input buses. In this chapter we design the
ALU our machine will use: we fix its contract, choose and encode eight
operations, define the four flags before writing a line of Verilog,
build it the \emph{compute-all-then-mux} way, prove it with two
self-checking benches (one directed, one exhaustive over
131{,}072 cases), and map it onto the Tang Nano 9K's switches and
LEDs. By the end you own the first real organ of the C1 capstone.

% from-zero:begin
\begin{rulebox}[Before you start]
Chapters~8--9 and~11--12. \emph{Recap:} an ALU is the part of a CPU that adds, subtracts and compares. It is the combinational block of Chapter~9 with a control input.
\end{rulebox}
\begin{wins}[After this chapter you can]
\begin{itemize}
\item specify an ALU's operations and flags before writing code;
\item implement an 8-bit ALU in Verilog;
\item test it with vectors that check the flags as well as the result;
\item run it on a board.
\end{itemize}
\end{wins}
\begin{principle}[How to think about this]
An ALU is a box that computes every operation at once and then picks the one you asked for. Decide the contract first, including what each flag means, then write the code. Testing the flags as carefully as the result is where most ALU bugs are caught.
\end{principle}
% from-zero:end

\section{What an ALU is, and the contract}\label{sec:contract}

\index{ALU|see{arithmetic-logic unit}}\index{arithmetic-logic unit}
An \emph{arithmetic-logic unit} (ALU) is the part of a processor that
combines two operand words into one result word, under the control of
an operation select field, and reports a few one-bit facts about the
answer. It is \emph{combinational}: given inputs, it answers
immediately. It holds no state --- the registers around it (ch15)
decide what to keep.

\begin{figure}[htbp]
\centering
\begin{tikzpicture}[font=\small, >=latex]
% the contract: a black box with named ports
\draw[thick, ink, fill=graybg] (0,0) rectangle (3.4,4.6);
\node[font=\bfseries\large] at (1.7,2.3) {alu8};
\node[font=\scriptsize, text=slate] at (1.7,1.9) {8-bit ALU};
% inputs (left)
\draw[->, thick] (-1.6,3.9) -- (0,3.9) node[midway, above, font=\scriptsize]{\texttt{i\_a}/8};
\draw[->, thick] (-1.6,3.1) -- (0,3.1) node[midway, above, font=\scriptsize]{\texttt{i\_b}/8};
\draw[->, thick] (-1.6,1.1) -- (0,1.1) node[midway, above, font=\scriptsize]{\texttt{i\_op}/3};
\node[left, font=\scriptsize, text=slate] at (-1.6,3.9) {operand A};
\node[left, font=\scriptsize, text=slate] at (-1.6,3.1) {operand B};
\node[left, font=\scriptsize, text=slate] at (-1.6,1.1) {which operation};
% outputs (right)
\draw[->, thick] (3.4,3.5) -- (5.0,3.5) node[midway, above, font=\scriptsize]{\texttt{o\_y}/8};
\node[right, font=\scriptsize, text=slate] at (5.0,3.5) {result};
\foreach \y/\n/\lbl in {2.5/z/zero, 1.9/n/negative, 1.3/c/carry, 0.7/v/overflow}{
  \draw[->, thick, accent] (3.4,\y) -- (5.0,\y);
  \node[right, font=\scriptsize] at (5.0,\y) {\texttt{o\_\n} = \lbl};}
\node[font=\scriptsize, text=slate, align=center] at (1.7,-0.5) {no clock, no memory:\\
\emph{combinational} --- answers now, not later};
\end{tikzpicture}

\caption{The contract as a black box. Two 8-bit operands and a 3-bit
operation field go in; an 8-bit result and four flags come out. Notice
what is absent: no clock and no reset --- a combinational block has
nothing to remember. Regenerate: \ttfamily figures/ch14/fig\_block.tex.}
\label{fig:block14}
\end{figure}

Writing the contract down first --- before any Verilog --- is the
single habit that separates a machine that works from a box of parts.
Ours:

\begin{description}
\item[\texttt{i\_a[7:0]}, \texttt{i\_b[7:0]}] the operands. Both are
always read; nothing is optional.
\item[\texttt{i\_op[2:0]}] the operation field. Three bits, eight
codes, every code assigned.
\item[\texttt{o\_y[7:0]}] the result.
\item[\texttt{o\_z}, \texttt{o\_n}, \texttt{o\_c}, \texttt{o\_v}] the
\emph{flags}\index{flag}: single-bit facts about the last result ---
zero, negative, carry, overflow. Defined precisely in
\S\ref{sec:flags14}.
\end{description}

That is the whole interface. Everything else --- how many adders, how
the mux is wired --- is our private business behind the port list.

\section{Choosing the operations}\label{sec:opset14}

What should the ALU compute? The honest answer is: whatever the
instruction set needs, plus the cheap extras that fall out of the same
gates. Our ISA\index{ISA} (finalised in ch16) needs add, subtract and
bitwise-and for its arithmetic instructions; comparisons fall out of
subtraction for free. The full set, encoded in three bits:

\begin{figure}[htbp]
\centering
\begin{tikzpicture}[font=\small]
\node[anchor=north west, inner sep=0] at (0,0) {%
\begin{tabular}{llll}
\toprule
\texttt{i\_op} & name & result \texttt{o\_y} & flags that can move\\
\midrule
\texttt{000} & \texttt{ADD} & $a+b$ & Z N \textbf{C V}\\
\texttt{001} & \texttt{SUB} & $a-b$ & Z N \textbf{C V}\\
\texttt{010} & \texttt{AND} & $a\mathbin{\&}b$ & Z N\\
\texttt{011} & \texttt{OR}  & $a\mathbin{|}b$ & Z N\\
\texttt{100} & \texttt{XOR} & $a\mathbin{\hat{}}b$ & Z N\\
\texttt{101} & \texttt{SHL} & $a\ll1$ & Z N \textbf{C}\\
\texttt{110} & \texttt{SHR} & $a\gg1$ & Z N \textbf{C}\\
\texttt{111} & \texttt{CMP} & $a-b$ & Z N \textbf{C V}\\
\bottomrule
\end{tabular}};
\node[anchor=north west, text width=6.2cm, font=\scriptsize, text=slate]
  at (0,-5.6) {The encoding is a contract, not a preference: the ch16
  instruction set drives arithmetic through these three bits, so the
  table is written in stone for the rest of the book. \texttt{CMP}
  shares the subtracter's datapath but the difference it computes is
  thrown away --- the flags are the verdict.};
\end{tikzpicture}

\caption{The op encoding table. Notice that \texttt{SUB} and
\texttt{CMP} do the same subtraction --- compare \emph{is} a subtract
nobody writes back. Also notice which ops are allowed to move the
arithmetic flags: only add, subtract, compare and the shifts touch
\texttt{C}/\texttt{V}; the bitwise ops leave them at 0.
Regenerate: \ttfamily figures/ch14/fig\_optable.tex.}
\label{fig:optable14}
\end{figure}

Two encodings decisions worth stating out loud:

\begin{itemize}
\item The three arithmetic instructions our ISA will use ---
\texttt{ADD}, \texttt{SUB}, \texttt{AND} --- sit at codes
\texttt{000}--\texttt{010}. When ch16's control unit decodes an
instruction it will drive these same three bits: the table above is
the contract between this chapter and that one.
\item Each of our eight ops has a direct ancestor in a real ISA's
arithmetic: \texttt{ADD}$\to$\texttt{add}, \texttt{SUB}$\to$\texttt{sub},
\texttt{AND/OR/XOR}$\to$\texttt{and/or/xor},
\texttt{SHL}$\to$\texttt{sll}, \texttt{SHR}$\to$\texttt{srl},
\texttt{CMP}$\to$a \texttt{sub} whose result is discarded. That is the
RISC-V (RV32I) mapping our ISA inherits in ch16 --- our toy machine
really is a small copy of the real thing.
\end{itemize}

\begin{exercise}[easy]
Why is \texttt{CMP} a separate code at all, when \texttt{SUB} already
computes $a-b$ and sets the same flags? What does the CPU actually do
differently for the two?
\end{exercise}

\section{Flags: semantics before code}\label{sec:flags14}

\index{flag!zero}\index{flag!negative}\index{flag!carry}\index{flag!overflow}
The flags are the ALU's way of telling the control unit \emph{what
kind} of answer just came out. Get their meanings on paper before
coding, or you will find yourself reading Verilog to discover what you
meant --- backwards.

\begin{definition}[The four flags]
Let $r$ be the 8-bit result of an operation. Then:
\begin{description}
\item[Z (zero)] $Z=1$ iff $r=\texttt{0x00}$.
\item[N (negative)] $N=r_7$: the sign bit, when $r$ is read as two's
complement.
\item[C (carry)] for \texttt{ADD}: the carry out of bit~7 --- unsigned
overflow. For \texttt{SUB}/\texttt{CMP}: the carry out of
$a+\bar b+1$, i.e.\ \emph{no borrow} --- $C=1$ means $a\geq b$
unsigned. For shifts: the bit that fell off the end. For logic ops: 0.
\item[V (overflow)] for \texttt{ADD}: $V=1$ iff $a_7=b_7$ and
$r_7\neq a_7$ (same-sign inputs, wrong-sign result). For
\texttt{SUB}/\texttt{CMP}: $V=1$ iff $a_7\neq b_7$ and $r_7\neq a_7$
(different-sign inputs, result's sign moved away from $a$'s). For all
other ops: 0.
\end{description}
\end{definition}

Read the V rule slowly --- it is the one flag beginners misimplement.
Signed overflow can \emph{only} happen when the true mathematical
answer does not fit in eight bits: adding two positives and getting a
negative, or subtracting a negative from a positive and getting a
negative. That is exactly ``inputs' sign bits agree (for add) /
disagree (for sub), and the result's sign bit is not what the first
operand's was.''

\begin{example}[Deriving the flag table, not asserting it]
Take \texttt{0x7F + 0x01}: $127+1=128$. In hex the result is
\texttt{0x80} --- which, as two's complement, is $-128$. Is that
overflow? Apply the rule: $a_7=0$, $b_7=0$ (same signs), $r_7=1\neq
a_7$ $\Rightarrow$ $V=1$. Carry out of the byte? $128<256$, so $C=0$.
Now \texttt{0xFF + 0x01}: $255+1=256$, which is \texttt{0x00} plus
carry $\Rightarrow$ $C=1$, $Z=1$. Overflow? $a_7=1$, $b_7=0$: signs
differ, so $V=0$ by definition --- and indeed $-1+1=0$ was perfectly
representable. The flags are not feelings; they are a lookup of the
definition.
\end{example}

\begin{checkpoint}
Before reading on: fill in Z N C V by hand for \texttt{0x80 + 0x80},
\texttt{0x00 - 0x01} (as \texttt{SUB}) and \texttt{0x80 - 0x01}.
Compare with the simulation rows in \S\ref{sec:tb14}.
\end{checkpoint}

\section{The design: compute all answers, then pick one}
\label{sec:design14}

The temptation is to write an ALU as an \texttt{if}-ladder that
``does'' add, or subtract, etc. Resist it: in hardware there is no
``only compute what you need.'' The wires exist, so they all answer,
all the time. The clean structure --- and the one synthesis rewards
--- is in Figure~\ref{fig:structure14}.

\begin{figure}[htbp]
\centering
\begin{tikzpicture}[font=\scriptsize, >=latex]
% operands fan out to every path at once
\node[left, font=\small] at (-0.4,4.6) {\texttt{i\_a}/8};
\node[left, font=\small] at (-0.4,4.0) {\texttt{i\_b}/8};
\draw[thick] (0,4.6) -- (0.6,4.6);
\draw[thick] (0,4.0) -- (0.6,4.0);
\draw[thick] (0.6,4.6) -- (0.6,0.6) -- (1.0,0.6);
\draw[thick] (0.6,4.0) -- (0.8,4.0) -- (0.8,0.9) -- (1.0,0.9);
% seven paths, always computing
\foreach \y/\nm/\det in {4.4/{adder}/{\texttt{a+b} \textrightarrow\ 9b},3.7/{subtracter}/{\texttt{a+\textasciitilde b+1} \textrightarrow\ 9b},3.0/{AND}/{\texttt{a\&b}},2.3/{OR}/{\texttt{a|b}},1.6/{XOR}/{\texttt{a\char`^b}},0.9/{SHL}/{\texttt{\{a[6:0],0\}}},0.2/{SHR}/{\texttt{\{0,a[7:1]\}}}}{
  \draw[thick, fill=graybg] (1.0,\y-0.22) rectangle (2.9,\y+0.22);
  \node at (1.95,\y) {\det};
}
% the big mux
\draw[thick, fill=accentsoft] (3.6,0.0) -- (4.0,2.0) -- (4.0,4.6) -- (3.6,4.8) -- cycle;
\node[rotate=90] at (3.85,2.4) {8$\times$8 mux};
\foreach \y in {4.4,3.7,3.0,2.3,1.6,0.9,0.2}{
  \draw[->, thick] (2.9,\y) -- (3.75,\y*0.94+0.15);
}
\draw[->, thick] (4.0,2.4) -- (5.4,2.4) node[right, font=\small]{\texttt{o\_y}/8};
% op field selects
\draw[->, thick, accent] (2.4,-0.7) -- (2.4,1.4) -- (3.65,1.4);
\node[below, font=\small] at (2.4,-0.7) {\texttt{i\_op}/3};
% flags ride alongside
\draw[->, thick, accent] (4.0,1.2) -- (4.6,1.2) -- (4.6,0.4) -- (5.4,0.4);
\node[right, font=\scriptsize] at (5.4,0.4) {\texttt{o\_z o\_n o\_c o\_v}};
\node[font=\scriptsize, text=slate, align=center] at (1.9,-1.3) {every path answers, always ---\\
the mux only picks};
\end{tikzpicture}

\caption{Compute-all-then-mux. The operands fan out to seven datapaths
that compute in parallel and permanently; the op field drives one big
mux that picks a result. The flag logic sits beside the mux, not
inside it. This is exactly how the classic teaching ALUs (and ours,
Listing~\ref{lst:alu8}) are built.}
\label{fig:structure14}
\end{figure}

The subtraction path deserves its own line of algebra. Two's
complement says $-b=\bar b+1$, so
\[
a-b \;=\; a+\bar b+1,
\]
and a 9-bit adder computing $a+\bar b+1$ produces both the difference
(bits 7--0) and, in bit~8, the answer to ``did we borrow?'' ---
inverted: carry-out set means \emph{no} borrow, the ARM convention we
adopted in \S\ref{sec:flags14}. One adder, minus sign, flags: three
birds.

\begin{example}[Subtract proven, not assumed]
Run it on \texttt{0x03 - 0x05}: $\bar b=\texttt{0xFA}$, so
$a+\bar b+1 = 3+250+1 = 254 = \texttt{0xFE}$, which is $-2$ in two's
complement --- correct. Bit~8? $3+250+1=254<256$, so no carry:
$C=0$, the borrow happened, as expected for $3<5$. Now
\texttt{0x05 - 0x03}: $5+252+1=258=\texttt{0x102}$: result
\texttt{0x02}, $C=1$, no borrow. The flag column of
Figure~\ref{fig:optable14} was earned, not guessed.
\end{example}

\begin{figure}[htbp]
\centering
\begin{tikzpicture}[font=\scriptsize, >=latex]
% selected result feeds Z and N unconditionally
\draw[thick, fill=graybg] (0,0) rectangle (2.2,1.1);
\node[align=center] at (1.1,0.55) {selected\\result \texttt{y}/8};
\draw[->, thick] (2.2,0.9) -- (3.6,0.9) -- (3.6,0.55) -- (4.4,0.55);
\draw[thick, fill=goodbg] (4.4,0.15) rectangle (6.4,0.95);
\node[align=center] at (5.4,0.55) {\texttt{y[7]}\\sign bit};
\draw[->, thick, accent] (6.4,0.55) -- (7.4,0.55) node[right]{\texttt{o\_n}};
\draw[->, thick] (2.2,0.25) -- (3.0,0.25) -- (3.0,-0.5) -- (3.8,-0.5);
\draw[thick, fill=goodbg] (3.8,-0.9) rectangle (6.0,-0.1);
\node[align=center] at (4.9,-0.5) {8-input NOR\\(\texttt{y == 0})};
\draw[->, thick, accent] (6.0,-0.5) -- (7.4,-0.5) node[right]{\texttt{o\_z}};
% carry: mux of per-path candidates
\draw[thick, fill=amberbg] (0.6,-2.4) rectangle (3.2,-1.2);
\node[align=center] at (1.9,-1.5) {\texttt{add[8]}};
\node[align=center] at (1.9,-1.9) {\texttt{sub[8]}};
\node[align=center] at (1.9,-2.3) {\texttt{a[7]}, \texttt{a[0]}, \texttt{0}};
\draw[thick, fill=graybg] (4.0,-2.5) rectangle (6.2,-1.3);
\node[align=center] at (5.1,-1.9) {mux on\\\texttt{i\_op}};
\draw[->, thick] (3.2,-1.8) -- (4.0,-1.8);
\draw[->, thick, accent] (6.2,-1.9) -- (7.4,-1.9) node[right]{\texttt{o\_c}};
% overflow: arithmetic paths only
\draw[thick, fill=amberbg] (0.6,-3.9) rectangle (3.2,-2.9);
\node[align=center] at (1.9,-3.2) {\texttt{v\_add}: same signs,\\new sign differs};
\node[align=center] at (1.9,-3.65) {\texttt{v\_sub}: signs differ,\\new sign $\neq$ \texttt{a[7]}};
\draw[thick, fill=graybg] (4.0,-3.95) rectangle (6.2,-3.0);
\node[align=center] at (5.1,-3.45) {mux on\\\texttt{i\_op}};
\draw[->, thick] (3.2,-3.4) -- (4.0,-3.4);
\draw[->, thick, accent] (6.2,-3.45) -- (7.4,-3.45) node[right]{\texttt{o\_v}};
\draw[->, dashed, slate] (5.1,-1.3) -- (5.1,-0.2) -- (1.1,-0.2) -- (1.1,0);
\draw[->, dashed, slate] (1.1,-1.2) -- (1.1,-0.2);
\end{tikzpicture}

\caption{Where the flags come from. $Z$ and $N$ are unconditional ---
they describe whatever result the mux chose. $C$ and $V$ are muxed on
\texttt{i\_op} alongside the result: each arithmetic path supplies its
own candidate, logic ops supply 0, and shifts supply the ejected bit.}
\label{fig:flags14}
\end{figure}

\section{The ALU in Verilog}\label{sec:rtl14}

Listing~\ref{lst:alu8} is the whole module. Read stage~1 as ``the
parallel answers'' and the \texttt{case} as the mux.

\begin{lstlisting}[language=Verilog, caption={\ttfamily code/ch14\_alu8.v
--- the ALU itself. Every path computes always; the \texttt{case} only
selects. \texttt{SUB} and \texttt{CMP} share one subtracter.},
label={lst:alu8}]
module alu8 (
  input  wire [7:0] i_a,
  input  wire [7:0] i_b,
  input  wire [2:0] i_op,
  output reg  [7:0] o_y,
  output reg        o_z,
  output reg        o_n,
  output reg        o_c,
  output reg        o_v
);
  localparam OP_ADD = 3'b000, OP_SUB = 3'b001, OP_AND = 3'b010,
             OP_OR  = 3'b011, OP_XOR = 3'b100, OP_SHL = 3'b101,
             OP_SHR = 3'b110, OP_CMP = 3'b111;

  // stage 1: every answer computed in parallel, always
  wire [8:0] add = {1'b0, i_a} + {1'b0, i_b};
  wire [8:0] sub = {1'b0, i_a} + {1'b0, ~i_b} + 9'd1;
  wire [7:0] y_and = i_a & i_b;
  wire [7:0] y_or  = i_a | i_b;
  wire [7:0] y_xor = i_a ^ i_b;
  wire [7:0] y_shl = {i_a[6:0], 1'b0};
  wire [7:0] y_shr = {1'b0, i_a[7:1]};

  wire c_add = add[8];
  wire v_add = (i_a[7] == i_b[7]) && (add[7] != i_a[7]);
  wire c_sub = sub[8];
  wire v_sub = (i_a[7] != i_b[7]) && (sub[7] != i_a[7]);

  // stage 2: the op field picks one result and one {C,V} pair
  always @(*) begin
    case (i_op)
      OP_ADD:         begin o_y = add[7:0]; o_c = c_add;  o_v = v_add; end
      OP_SUB, OP_CMP: begin o_y = sub[7:0]; o_c = c_sub;  o_v = v_sub; end
      OP_AND:         begin o_y = y_and;    o_c = 1'b0;   o_v = 1'b0;  end
      OP_OR:          begin o_y = y_or;     o_c = 1'b0;   o_v = 1'b0;  end
      OP_XOR:         begin o_y = y_xor;    o_c = 1'b0;   o_v = 1'b0;  end
      OP_SHL:         begin o_y = y_shl;    o_c = i_a[7]; o_v = 1'b0;  end
      OP_SHR:         begin o_y = y_shr;    o_c = i_a[0]; o_v = 1'b0;  end
      default:        begin o_y = 8'h00;    o_c = 1'b0;   o_v = 1'b0;  end
    endcase
    o_z = (o_y == 8'h00);
    o_n = o_y[7];
  end
endmodule
\end{lstlisting}

Three details to notice. The complement is written
\texttt{\{1'b0,\textasciitilde i\_b\}} on purpose: inside a 9-bit
expression, bare \texttt{\textasciitilde i\_b} would be widened to
9 bits \emph{first} and complemented \emph{second}, giving
\texttt{9'h1FE} instead of \texttt{9'h0FE} --- the context-width trap
in Common mistakes. And CMP drives the difference onto
\texttt{o\_y} just like SUB: the CPU's write-enable (ch15) is what
turns a subtraction into a comparison, not anything inside the ALU.
Finally, \texttt{o\_z}/\texttt{o\_n} sit \emph{after} the case: they
describe the selected result whoever produced it.

\section{The testbench: vectors plus flags}\label{sec:tb14}

A directed bench can't prove much by staring at outputs --- so ours
reads a vector file where each row carries the \emph{expected} result
and all four expected flags, generated by an independent Python
golden model (\texttt{code/ch14\_alu8\_vec\_gen.py}). The bench
applies each vector, waits, then compares all 12 output bits with
\texttt{!==} so an \texttt{x} anywhere counts as failure:

\begin{lstlisting}[caption={Real run of \ttfamily ch14\_alu8\_tb.v ---
the directed vectors. Each row prints what was applied and what came
back; the verdict counts checks, not hopes.}]
t=50 op=000 a=7f b=01 -> y=80 z=0 n=1 c=0 v=1
t=60 op=000 a=80 b=80 -> y=00 z=1 n=0 c=1 v=1
t=70 op=000 a=ff b=01 -> y=00 z=1 n=0 c=1 v=0
t=80 op=000 a=40 b=40 -> y=80 z=0 n=1 c=0 v=1
t=90 op=000 a=c0 b=40 -> y=00 z=1 n=0 c=1 v=0
t=100 op=001 a=05 b=03 -> y=02 z=0 n=0 c=1 v=0
t=110 op=001 a=03 b=05 -> y=fe z=0 n=1 c=0 v=0
t=120 op=001 a=00 b=00 -> y=00 z=1 n=0 c=1 v=0
t=130 op=001 a=80 b=01 -> y=7f z=0 n=0 c=1 v=1
t=140 op=001 a=7f b=ff -> y=80 z=0 n=1 c=0 v=1
t=150 op=001 a=0a b=0a -> y=00 z=1 n=0 c=1 v=0
...
t=640 op=000 a=00 b=00 -> y=00 z=1 n=0 c=0 v=0
PASS: all 64 vectors match (result + ZN CV flags)
\end{lstlisting}

Check t=50 against your hand derivation from \S\ref{sec:flags14}:
\texttt{0x7F + 0x01 = 0x80} gives $N=1$, $V=1$, $C=0$ --- the waveform
in Figure~\ref{fig:wave14} shows the same transition. And t=120 is the
row most people get wrong on paper: $0-0$ leaves $Z=1$ \emph{and}
$C=1$, because $0+\texttt{0xFF}+1$ carried out --- no borrow needed.

\begin{figure}[htbp]
\centering
\begin{tikzpicture}
\begin{axis}[width=0.85\linewidth, height=4.2cm, xlabel={$t$ (ns)},
  xmin=0, xmax=175, ymin=-0.3, ymax=1.4, clip=false,
  ytick={0,1}, yticklabels={0,1}, axis lines=left]
\addplot[const plot, thick, blue] table[x=t,y=v] {figures/ch13/ch13_blink.dat};
\node[anchor=east, font=\small] at (axis cs:0,0.5) {\sf led};
\node[font=\scriptsize, text=gray] at (axis cs:115,1.15) {bit 3 high for 8 clocks};
\end{axis}
\end{tikzpicture}

\caption{The flag edge cases on tape: the real \texttt{alu8\_tb} VCD
rendered. Watch t=50--60 (ps$\times10^3$): \texttt{0x80+0x80} sets
$Z$, $C$ \emph{and} $V$ at once; one vector later \texttt{0xFF+0x01}
keeps $Z$ and $C$ but drops $V$ --- same all-zero result, different
flags. That distinction is the entire point of having two flags.}
\label{fig:wave14}
\end{figure}

\begin{example}[Directed vectors are not enough --- so we sweep all of them]
Sixty-four hand-picked rows cover the cases we \emph{imagined}.
\texttt{ch14\_alu8\_exh\_tb.v} does the other thing: it nests two
256-loops and drives \emph{every} operand pair for ADD, then every
pair for SUB, scoring against a golden model written from the flag
definitions --- a different description of the same truth, so a bug in
one can't hide in the other:

\begin{lstlisting}[caption={Real run of \ttfamily ch14\_alu8\_exh\_tb.v
--- exhaustive proof for the arithmetic paths.}]
ADD done: 65536 checks, 0 errors
PASS: 131072 exhaustive vectors clean
\end{lstlisting}

That is stronger evidence than any amount of directed testing, and it
took 0.4 seconds. Coverage this cheap is a bargain --- take it when
the input space is small enough.
\end{example}

\begin{example}[What it costs in silicon]
Strategy: run \texttt{alu8} through \texttt{synth\_gowin} like ch13
and read the \texttt{stat} block:

\begin{lstlisting}[caption={\ttfamily synth\_gowin stat for
\texttt{alu8} --- real output, trimmed.}]
=== alu8 ===
      239 wires      345 wire bits        8 ports
      263 cells
       18   ALU
       53   LUT1        6   LUT2       22   LUT3       43   LUT4
       54   MUX2_LUT5  24   MUX2_LUT6   8   MUX2_LUT7   2   MUX2_LUT8
       19   IBUF       12   OBUF         1   GND         1   VCC
\end{lstlisting}

Line by line: \textbf{18 ALU} cells are the two 8-bit adders mapped
onto the fast-carry tiles (the mux cells can't do arithmetic).
\textbf{124 LUTs} cover the logic ops, shifts and flag expressions.
\textbf{88 wide mux cells} (\texttt{MUX2\_LUT5/6/7/8}) are the price
of compute-all-then-mux --- that big 8-way select is where most of the
area went, exactly as Figure~\ref{fig:structure14} predicts.
\textbf{19 IBUF + 12 OBUF} = 19 inputs + 12 outputs: the port list,
accounted for. Nothing in the report is magic once the structure is in
your head.
\end{example}

\section{Onto the board}\label{sec:board14}

Simulation proved the ALU; the board makes it tangible. The demo
--- \texttt{ch14\_top\_alu9k.v} and \texttt{ch14\_alu9k.cst}, both in
\texttt{code/} --- needs no clock at all --- a nice mirror of the ALU's own
combinational nature. Two DIP-8 banks on the headers drive
\texttt{i\_a} and \texttt{i\_b} live; three more switch positions drive
\texttt{i\_op}; the six onboard LEDs show the answer, and the S1
button swaps what they display.

\begin{figure}[htbp]
\centering
\begin{tikzpicture}[font=\scriptsize, >=latex]
% breadboard banks (left)
\draw[thick, fill=amberbg] (0,3.2) rectangle (3.0,4.6);
\node[align=center] at (1.5,4.15) {DIP-8 bank\\\texttt{i\_a[7:0]}};
\node[align=center, text=slate] at (1.5,3.5) {one leg to GND\\open leg = 1 (pull-up)};
\draw[thick, fill=amberbg] (0,1.5) rectangle (3.0,2.9);
\node[align=center] at (1.5,2.45) {DIP-8 bank\\\texttt{i\_b[7:0]}};
\node[align=center, text=slate] at (1.5,1.8) {same trick on J5};
\draw[thick, fill=amberbg] (0,0.0) rectangle (3.0,1.2);
\node[align=center] at (1.5,0.6) {DIP-3 / jumpers\\\texttt{i\_op[2:0]}};
% board box (right)
\draw[thick, ink] (4.2,-0.3) rectangle (8.4,5.0);
\node[font=\small\bfseries] at (6.3,4.65) {Tang Nano 9K};
\draw[thick, fill=graybg] (4.6,3.4) rectangle (6.2,4.2);
\node[align=center] at (5.4,3.8) {J6 header\\79--86, 76, 77, 48};
\draw[thick, fill=graybg] (4.6,1.9) rectangle (6.2,2.7);
\node[align=center] at (5.4,2.3) {J5 header\\25--30, 33, 34};
\draw[thick, fill=accentsoft] (6.6,2.2) rectangle (8.0,3.6);
\node[align=center] at (7.3,2.9) {\texttt{alu8}\\(+ mux)};
\draw[thick, fill=goodbg] (6.6,0.8) rectangle (8.0,1.9);
\node[align=center] at (7.3,1.5) {6 LEDs\\active-low};
\node[align=center, text=slate] at (7.3,1.05) {10,11,13,14,15,16};
\draw[thick, fill=amberbg] (6.6,-0.1) rectangle (8.0,0.5);
\node[align=center] at (7.3,0.2) {S1 button (pin 3)};
% wires
\draw[->, thick] (3.0,3.9) -- (4.6,3.9);
\draw[->, thick] (3.0,2.3) -- (4.6,2.3);
\draw[->, thick] (3.0,0.6) -- (4.0,0.6) -- (4.0,3.5) -- (4.6,3.5);
\draw[->, thick] (6.2,3.8) -- (6.6,3.4);
\draw[->, thick] (6.2,2.3) -- (6.6,2.6);
\draw[->, thick] (7.3,2.2) -- (7.3,1.9);
\draw[->, thick, accent] (7.3,0.5) -- (7.3,0.8);
\node[font=\scriptsize, text=slate, align=center] at (5.4,-0.7) {pin numbers: schematic rev 3674\\
--- verify against your board};
\end{tikzpicture}

\caption{The wiring plan. Switch legs go to GND so an open contact
reads 1 through the pin's pull-up; header pins follow the rev-3674
schematic (\emph{check your own rev before soldering}). Six LEDs can't
show 8 result bits and 4 flags at once, so S1 multiplexes the display:
released = \texttt{y[5:0]}, held = \texttt{\{Z,N,C,V,y[7:6]\}}.
Regenerate: \ttfamily figures/ch14/fig\_board.tex.}
\label{fig:board14}
\end{figure}

One hardware trap the schematic warns about: the LEDs are wired
\emph{active-low} --- a pin sinks current to light its LED, so the
wrapper inverts the display lines (\texttt{o\_led = \textasciitilde
disp}). If your first run shows ``the bits I expected, upside-down,''
that is why.

\begin{example}[Predicting the board before you flash]
Set \texttt{i\_a=0x0F}, \texttt{i\_b=0x01}, \texttt{i\_op=000}
(\texttt{ADD}). The mux picks the adder path:
$y=\texttt{0x10}$, $Z=0$, $N=0$, $C=0$, $V=0$. S1 released, the six
LEDs show \texttt{y[5:0]=0x10} --- only the bit-4 LED lit. Hold S1 and
the same six LEDs show
\texttt{\{Z,N,C,V,y[7:6]\}=0b000000} --- all dark, because the result
carried no flags and its top bits are zero. Now switch
\texttt{i\_b} to \texttt{0x7F}: $y=\texttt{0x8E}$, and holding S1
lights the N-flag bit and the \texttt{y[7]} bit --- the two's
complement wrap you derived in \S\ref{sec:flags14}, on physical
hardware. Working the demo backwards like this --- predict bits,
then read LEDs --- is how you distinguish a design bug from a
wiring bug.
\end{example}

\textbf{Honest marking.} The vector test, the exhaustive sweep and
the synthesis stat in this chapter were run and captured on this
machine. The wiring on Figure~\ref{fig:board14} is drawn from the
Sipeed rev-3674 schematic, but no board was attached here --- the
build commands (\texttt{synth\_gowin} $\to$
\texttt{nextpnr-himbaechel} $\to$ \texttt{gowin\_pack} $\to$
\texttt{openFPGALoader}) are the ones you would run, unchanged from
ch13's flow.

\section{Common mistakes}\label{sec:cm14}

\begin{tripwire}[Treating carry and overflow as the same flag]
\texttt{0xFF+0x01} and \texttt{0x7F+0x01} both make sense only with
two flags: the first wraps the \emph{unsigned} byte ($C=1$, $V=0$),
the second breaks the \emph{signed} byte ($V=1$, $C=0$). If your V
logic ever mentions the carry bit, it is wrong --- V is computed from
sign bits alone.
\end{tripwire}
\begin{tripwire}[Reading the C flag on SUB as ``borrow'']
We use the ARM convention: $C=1$ after SUB means \emph{no} borrow
($a\geq b$ unsigned). x86 does it the other way and neither is wrong
--- but mixing the two inside one machine is. Whatever you pick,
write it in the flag table and test it: vector \texttt{00-00} is the
smallest case where the conventions visibly differ.
\end{tripwire}
\begin{tripwire}[The context-width complement]
In a 9-bit expression, \texttt{\textasciitilde i\_b} is widened
\emph{then} complemented: \texttt{\textasciitilde 8'h01} inside
\texttt{a + \textasciitilde b + 9'd1} gives \texttt{9'h1FE}, not
\texttt{9'h0FE} --- an off-by-256 bug our own golden model committed
and the exhaustive sweep caught instantly. Complement at the operand's
width, then extend: \texttt{\{1'b0,\textasciitilde i\_b\}}.
\end{tripwire}
\begin{tripwire}[An op code with no owner]
Every one of the eight 3-bit codes is assigned --- but only because we
chose that. Add a ninth op someday and the code space is full;
silently aliasing two ops to one code is how off-by-one ISA bugs are
born. The \texttt{default} arm returning 0 is the seatbelt; the
exercise on op \texttt{111} asks you to think about what the ALU owes
a caller who feeds it a code you never meant.
\end{tripwire}
\begin{tripwire}[Expecting flags to mean something during a logic op]
\texttt{AND}/\texttt{OR}/\texttt{XOR} force $C=V=0$ by definition ---
a flag that sometimes means ``carry'' and sometimes means nothing is a
flag nobody can rely on. Whatever your table says for a row, the bench
must check that row too.
\end{tripwire}
\begin{tripwire}[Forgetting the LEDs are active-low]
On the Tang Nano 9K a \texttt{1} on the LED pin means dark. If the
demo shows the complement of your expected pattern, the design is
fine --- the polarity wasn't.
\end{tripwire}

\begin{exercise}[easy]
By hand: \texttt{SUB} with $a=\texttt{0x30}$, $b=\texttt{0x10}$.
Give \texttt{o\_y} and all four flags, then check against the table.
\end{exercise}
\begin{exercise}[medium]
\texttt{SHL} of \texttt{0xC5}: result and flags? Now two \texttt{SHL}s
in a row --- what does the second C flag report, and why is that a
useful behaviour for a 16-bit shift built from two 8-bit ops?
\end{exercise}
\begin{exercise}[medium]
Op field \texttt{111} is \texttt{CMP} in our table. Trace what the
ALU drives on \texttt{o\_y} during a compare, then decide: should the
CPU's datapath write that result back to a register? Write the
one-line case change you would make if the answer were ``no, force
\texttt{o\_y=0}'' --- and argue whether it matters which you pick.
\end{exercise}
\begin{exercise}[hard]
Prove from the stat report that swapping \texttt{SUB} onto a shared
adder saved area: estimate what a third 8-bit adder would add in
\texttt{ALU} cells, and what the \texttt{case} sharing
\texttt{OP\_SUB, OP\_CMP} cost in mux cells instead.
\end{exercise}

\section{Summary}\label{sec:sum14}

\begin{center}
\begin{tabularx}{\textwidth}{l>{\raggedright\arraybackslash}X}
\toprule
\textbf{Thing} & \textbf{The rule of thumb}\\
\midrule
ALU contract & two operands + op field in; result + Z N C V out;
combinational, no clock\\
op encoding & 3 bits, 8 codes, all assigned --- the contract ch16's
ISA drives\\
Z N C V & zero = all-zero result; negative = sign bit; carry =
unsigned wrap or ejected bit; overflow = signed wrap from sign bits
alone\\
$a-b$ & $=a+\bar b+1$; bit 8 = no-borrow (ARM convention: $C=1$ means
$a\geq b$)\\
structure & compute every op always; a mux selects --- don't write
``only compute what you need''\\
CMP & a SUB whose result nobody writes back; flags are the answer\\
verification & directed vectors for edge cases; exhaustive sweep when
the space is small; golden model written from spec, not DUT\\
synthesis reads & 18 fast-carry ALU cells for the adders; ~88 mux
cells is the price of compute-all-then-mux\\
\bottomrule
\end{tabularx}
\end{center}

\section*{Additional exercises}
\addcontentsline{toc}{section}{Additional exercises}

\begin{exercise}[easy]
List the eight op codes and names from memory, and say which two share
one datapath.
\end{exercise}
\begin{exercise}[easy]
For \texttt{AND} of \texttt{0xF0} and \texttt{0x0F}: give
\texttt{o\_y}, Z, N, C, V.
\end{exercise}
\begin{exercise}[easy]
Which flag would a conditional-jump instruction read to implement
``jump if equal'' after a \texttt{CMP}? Which for ``jump if
unsigned-below''?
\end{exercise}
\begin{exercise}[easy]
\texttt{SHR} of \texttt{0x81}, then of the result again: show result,
C and Z after each step.
\end{exercise}
\begin{exercise}[medium]
Fill the full flag row for \texttt{ADD 0xFF 0xFF}, \texttt{SUB 0x00
0x80} and \texttt{CMP 0x80 0x80} without running the sim; then check
with the vector file format from \S\ref{sec:tb14}.
\end{exercise}
\begin{exercise}[medium]
The tb compares with \texttt{!==}. Explain what plain \texttt{!=}
would miss, and which flag bit is most likely to go \texttt{x} on a
buggy board.
\end{exercise}
\begin{exercise}[medium]
Design op \texttt{NOT} ($\bar a$) for this ALU. There is no free code
--- name two honest ways to fit it, and pick one.
\end{exercise}
\begin{exercise}[medium]
A teammate writes \texttt{wire [8:0] sub = i\_a - i\_b;}. Does it
produce the same bit-9 result and flags as our
$a+\bar b+1$ form? Answer carefully: same numeric result, but what
changes for the C flag, and does \texttt{iverilog} care?
\end{exercise}
\begin{exercise}[medium]
Sketch the flag-logic change needed if \texttt{SHL} should set V like
x86 does (V = sign change of the result). Is it consistent with our
flag table as written?
\end{exercise}
\begin{exercise}[hard]
Extend the vector file and bench with a ninth pseudo-op ``run all 8
ops on this (a,b) pair'' by driving \texttt{i\_op} through all codes
inside the vector loop. What new bug class does that coverage expose?
\end{exercise}
\begin{exercise}[hard]
Write the exhaustive sweep for all \emph{six} remaining ops (logic and
shifts) --- $256^2$ for logic, $256$ for each shift --- and argue why
exhaustive checking of bitwise ops is still worth 40 ms of CPU time.
\end{exercise}
\begin{exercise}[hard]
Our V rule for \texttt{SUB}: $a_7\neq b_7$ and $r_7\neq a_7$. Prove
it is equivalent to ``signed $a-b$ does not fit in 8 bits'' by
enumerating the four sign combinations of $a$ and $b$.
\end{exercise}
\begin{exercise}[hard]
In Figure~\ref{fig:structure14}, the OR/XOR/AND paths are single-gate
deep while the adder is the long pole. If this ALU fed a 27\,MHz
register, which path bounds the clock, and which structural change
(not lower clock) removes the bound?
\end{exercise}
\begin{exercise}[hard]
The mux cells (\texttt{MUX2\_LUT*}) are ~35\% of the cell count.
An alternative structure computes the result bit-serially over 8
clocks. Estimate the flop/LUT cost of the serial version and name the
two things it gains and the one thing it gives up.
\end{exercise}
\begin{exercise}[hard]
Write \texttt{alu8\_r}: same contract, but \texttt{i\_op} takes the
4-bit value of a ch16 ISA opcode instead (assume opcodes
\texttt{ADD=4'b0010}, \texttt{SUB=4'b0011}, \texttt{AND=4'b0100}). What
belongs in the ALU versus in the control unit's decoder? Argue both
sides.
\end{exercise}
\begin{exercise}[hard]
Plant a bug: change \texttt{o\_z} to \texttt{(o\_y == 8'h01)}. Which
of the 64 vectors fire, and which vector would you add to guarantee
the sabotage is caught on the first run?
\end{exercise}
\begin{exercise}[hard]
The exhaustive sweep treats ADD and SUB; design the analogous
``exhaustive'' argument for \texttt{SHL}/\texttt{SHR} that is not a
loop over inputs at all --- a one-paragraph proof by induction on the
bit position.
\end{exercise}
\begin{exercise}[hard]
Read \texttt{synth\_gowin}'s report again: the \texttt{MUX2\_LUT8}
count is 2. Why only 2, when the result mux is 8-way? Explain the mux
tree depth and which result bit(s) actually needed the 8-input cell.
\end{exercise}
\begin{exercise}[hard]
Define the flag contract for an \texttt{ADC} (add-with-carry) op
taking $a+b+C_{in}$: truth table for the new C, and one sentence on
why multi-byte arithmetic is impossible without it.
\end{exercise}

\section*{Hints and answers}
\addcontentsline{toc}{section}{Hints and answers}
\begin{description}
\item[CMP vs SUB (easy, inline)] Same datapath and flags; the
difference is whether the CPU writes \texttt{o\_y} anywhere. A
compare's contract is ``flags only.''
\item[Flags hand check (checkpoint)] \texttt{0x80+0x80}: $y=00$, $Z=1$,
$N=0$, $C=1$, $V=1$. \texttt{0x00-0x01}: $y=\texttt{FF}$, $Z=0$,
$N=1$, $C=0$ (borrow), $V=0$. \texttt{0x80-0x01}: $y=\texttt{7F}$,
$Z=0$, $N=0$, $C=1$, $V=1$.
\item[0x30$-$0x10 (easy)] $y=\texttt{0x20}$; $Z=0$, $N=0$, $C=1$
(no borrow), $V=0$.
\item[SHL 0xC5 (medium)] $y=\texttt{0x8A}$, $C=1$ (top bit ejected),
$N=1$, $Z=0$, $V=0$; second \texttt{SHL} gives \texttt{0x14}, $C=1$.
Chain the C out of one byte into the LSB of the next --- the classic
carry-chain shift.
\item[CMP's result (medium)] \texttt{o\_y}=$a-b$, live on the bus.
Whether it matters is ch15's write-enable question; forcing
\texttt{o\_y=0} on CMP is legal but wastes a signal that costs
nothing to leave connected.
\item[Area proof (hard)] A third adder $\approx$ +9 ALU cells and a
wider input mux; sharing SUB/CMP instead costs one extra arm in an
already-wide mux --- cheap, which is why real ALUs share.
\item[Eight codes (easy)] \texttt{ADD} \texttt{SUB} \texttt{AND} \texttt{OR} \texttt{XOR} \texttt{SHL} \texttt{SHR} \texttt{CMP};
SUB/CMP share the subtracter.
\item[AND 0xF0\&0x0F (easy)] $y=0$; $Z=1$, $N=0$, $C=0$, $V=0$.
\item[Jump-if-equal (easy)] $Z$ for equal; $C=0$ means below
(unsigned borrow) --- or $N$ for signed-negative tests.
\item[SHR chain (easy)] \texttt{0x81$\to$0x40} $C=1$, then
\texttt{0x40$\to$0x20} $C=0$; $Z=0$ both.
\item[Flag rows (medium)] \texttt{FF+FF}: $y=\texttt{FE}$, $Z=0$,
$N=1$, $C=1$, $V=0$. \texttt{00$-$80}: $y=\texttt{80}$, $Z=0$, $N=1$,
$C=0$, $V=1$. \texttt{80$-$80}: $y=0$, $Z=1$, $N=0$, $C=1$, $V=0$.
\item[\texttt{!==} vs \texttt{!=} (medium)] \texttt{x != 0} evaluates
to \texttt{x}, which an \texttt{if} treats as false --- unknowns skip
the alarm. Any flag computed from an un-driven operand is the likely
\texttt{x}.
\item[NOT op (medium)] Share an existing code (e.g.\ make
\texttt{SHL} a rotate and free one) or widen to 4 bits. Either is
honest; widening is the cleaner contract.
\item[i\_a$-$i\_b (medium)] Same 9-bit result, but Verilog's
\texttt{-} gives borrow-semantics in bit 8 differently across tools;
write the explicit $a+\bar b+1$ so the carry chain is visible.
\item[SHL V flag (medium)] $V=a_7\oplus r_7$ would be a rule change,
not a bug --- consistency means updating the flag table and the bench
together.
\item[All-ops vector (hard)] Drives every code on shared datapaths ---
catches ops that read stale per-path wires.
\item[Exhaustive logic (hard)] $256^2\times3$ logic vectors plus
$256\times2$ shift vectors still finish in milliseconds; bitwise ops
have no edge cases you can imagine better than all of them.
\item[V equivalence (hard)] Only $a\geq0,b<0$ with result $<0$, or
$a<0,b\geq0$ with result $\geq0$, is impossible in range --- the two
cases are exactly $a_7\neq b_7 \wedge r_7\neq a_7$.
\item[Critical path (hard)] The adder's carry chain; pipeline the
result (ch15) or accept it --- at 27\,MHz a 9-bit ripple adder is
nowhere near the bound anyway.
\item[Serial ALU (hard)] One adder + 8-clock sequencer: tiny LUT count
but +8 flops and an FSM; gains area, loses latency --- and a clock
domain the combinational version never needed.
\item[ISA-driven op (hard)] The 4$\to$3 map is decode, and decode is
the control unit's job; keep \texttt{alu8}'s 3-bit contract pure so
the ALU stays testable alone.
\item[Planted bug (hard)] The \texttt{00-00} SUB row and every
\texttt{y=0} row catch it; add \texttt{ADD 0x00 0x01} (expect $Z=0$)
as the always-false guard vector.
\item[Shift proof (hard)] Induction on bit $i$: bit $i$ of
\texttt{y\_shl} equals $a_{i-1}$ by construction; the flag is $a_7$;
no input space to enumerate.
\item[MUX depth (hard)] 8-wide mux = 3 levels of 2-input cells; only
the lowest-fanin bits bottom out at \texttt{LUT8} --- most bits
resolve at shallower cells.
\item[ADC contract (hard)] $C_{out}$ = carry of $a+b+C_{in}$
computed as a single 9-bit sum; without carry-in, byte $n$ can't see
byte $n-1$'s wrap.
\end{description}

\section*{Further reading}
\addcontentsline{toc}{section}{Further reading}

\begin{itemize}
\item Nisan \& Schocken, \emph{The Elements of Computing Systems}
(nand2tetris), Project~2 ---\newline
\url{nand2tetris.org/project02} --- build the Hack ALU in
HDL; same compute-all-then-mux spine as ours, different flag
vocabulary (\texttt{zr}, \texttt{ng} --- no C/V, worth asking why).
Intermediate. Verified 200, 2026-10-11.
\item Harris \& Harris, \emph{Digital Design and Computer
Architecture}, 2nd ed., \S5.2.4 (ALU) and ch.~7 (single-cycle
datapath) --- the textbook treatment our ch14--16 arc follows;
companion slides at \url{https://pages.hmc.edu/harris/ddca/}~\cite{ddca}.
Intermediate. Verified 200, 2026-10-11.
\item Lushay Labs, \emph{Our First CPU} (Tang Nano series, part~10).
\url{learn.lushaylabs.com/tang-nano-series} --- a working
8-bit CPU on this exact board; read the ALU portion now, the rest
becomes ch15/16 homework. Beginner--intermediate. Verified 200,
2026-10-11.
\item \emph{HDLBits} --- \url{https://hdlbits.01xz.net} ---
Circuits/Combinational ``Arithmetic Circuits'' drills the adder and
overflow idioms used here, with instant judge feedback. Beginner.
Verified 200, 2026-10-11.
\item ChipVerify, \emph{Verilog tutorial}.\newline
\url{chipverify.com/tutorials/verilog} --- the
\texttt{case}/mux and testbench idioms this chapter leans on, at
reference depth. Beginner. Verified 200, 2026-10-11.
\item Wikipedia, \emph{Arithmetic logic unit} and \emph{Two's
complement}.
\url{en.wikipedia.org/wiki/Arithmetic_logic_unit} and
\url{en.wikipedia.org/wiki/Two's_complement} --- for the
historical record and a second derivation of the flag rules.
Reference. Verified 200, 2026-10-11.
\end{itemize}

\begin{rulebox}[{Watch instead}]
\begin{enumerate}
\item \emph{Building the ALU} --- Ben Eater, 16:17,
\url{https://www.youtube.com/watch?v=S-3fXU3FZQc} --- the same block
built from 74-series chips on a breadboard; watch the flag lines
settle after each op, which is exactly our Figure~\ref{fig:wave14}.
Matches \S\ref{sec:contract}--\S\ref{sec:design14}. Opened and
confirmed 2026-10-11.
\item \emph{How Computers Calculate --- the ALU: Crash Course Computer
Science \#5} --- CrashCourse, 11:09,
\url{https://www.youtube.com/watch?v=1I5ZMmrOfnA} --- what an ALU is
and why flags exist; lighter, faster, good before the Verilog.
Matches \S\ref{sec:contract}--\S\ref{sec:flags14}. Opened and
confirmed 2026-10-11.
\end{enumerate}
\end{rulebox}

\chapternotes{The op set and compute-all-then-mux structure follow
the Hack ALU of nand2tetris project~2 and the ALU treatment in
Harris~\&~Harris~\cite{ddca}; the flag vocabulary (Z N C V) follows
the ARM convention so `carry' after a subtract means no-borrow. Board
pin numbers and the J5/J6 header map are from the Sipeed Tang Nano 9K
rev-3674 schematic~\cite{tangnano9k}; synthesis figures are a real
\texttt{yosys} 0.69 \texttt{synth\_gowin} run \cite{apicula}. The
RISC-V opcode mapping is seeded here and completed in ch16 per the
approved decisions.}
